\documentclass[11pt,a4paper]{article}

\usepackage[utf8]{inputenc}
\usepackage[T1]{fontenc}
\usepackage{amsmath,amssymb,amsthm}
\usepackage{mathtools}
\usepackage{bm}
\usepackage{graphicx}
\usepackage{booktabs}
\usepackage{hyperref}
\usepackage[margin=1in]{geometry}
\usepackage{algorithm}
\usepackage{algorithmic}
\usepackage{subcaption}
\usepackage{xcolor}
\usepackage{cleveref}

\newtheorem{definition}{Definition}
\newtheorem{proposition}{Proposition}
\newtheorem{theorem}{Theorem}
\newtheorem{remark}{Remark}

\newcommand{\field}{\bm{\phi}}
\newcommand{\evap}{\varepsilon}
\newcommand{\retain}{\alpha}
\newcommand{\R}{\mathbb{R}}
\newcommand{\E}{\mathbb{E}}
\newcommand{\x}{\bm{x}}
\newcommand{\q}{\bm{q}}
\newcommand{\K}{\bm{k}}
\newcommand{\vv}{\bm{v}}

\title{Dual-Channel Attention:\\
Critical Phase Transitions and Trail Fields\\
in Neural Sequence Processing}

\author{Research Notes}
\date{\today}

\begin{document}
\maketitle

\begin{abstract}
Systems poised at a critical phase transition---between quiescent and
runaway-active regimes---exhibit maximal dynamic range, information
transmission, and adaptability.  Neuronal avalanches in cortex and
pheromone trail formation in ant colonies both belong to the directed
percolation universality class, governed by the same three minimal
ingredients: \emph{amplification} (positive feedback),
\emph{injection} (activity-dependent deposition), and \emph{erosion}
(uniform decay)---subject to the constraint, demonstrated in a
companion lattice ant study, that erosion must be \emph{uncoupled}
from the field it erodes.  We ask whether importing these ingredients
and this constraint into transformer attention produces measurable
computational benefits.  We
propose \emph{dual-channel attention}, which decomposes the attention
mechanism into a \emph{structured channel} (content-based Q/K
matching) and a \emph{social channel}: a per-head vector recurrence
whose output modulates the key geometry of attention, reshaping
\emph{what} the model matches against rather than merely biasing
\emph{how much} attention each position receives.  The social channel
implements all three critical ingredients within the attention
computation itself, providing intrinsic sequential memory that is
architecturally distinct from content matching.  We ground this
design in the neurobiology of synaptic tagging, volume transmission,
and activity-dependent myelination---mechanisms that are functionally
stigmergic---and frame the architecture as an exploration of
criticality in learned sequence processing.
\end{abstract}


% ================================================================
\section{Introduction}
\label{sec:intro}

% ----------------------------------------------------------------
\subsection{Criticality in Learning Systems}

A growing body of evidence suggests that biological neural systems
operate near a \emph{critical phase transition}---the boundary between
a quiescent regime (activity dies out) and a runaway-active regime
(activity explodes).  At this critical point, systems exhibit maximal
dynamic range, maximal information transmission, and maximal
sensitivity to input \cite{beggs2003,shew2013,chialvo2010}.

Beggs and Plenz~\cite{beggs2003} showed that spontaneous activity in
cortical slices produces \emph{neuronal avalanches} whose size
distribution follows a power law with exponent $-3/2$---the signature
of the \emph{mean-field directed percolation} (DP) universality
class.  The branching parameter $\sigma$ converges to $\sim\!1$:
below 1, activity dies; above 1, it explodes; at exactly 1, it
persists at all scales.

This is not unique to neurons.  Munoz~\cite{munoz2018}, in a
comprehensive review in \emph{Reviews of Modern Physics}, shows that
the same critical scaling laws appear across living systems---from
gene regulatory networks to bird flocks to insect colony foraging.
The Janssen--Grassberger conjecture \cite{janssen1981,grassberger1982}
provides the theoretical bridge: \emph{any} system with (1)~a
continuous phase transition into an absorbing state, (2)~a
one-component order parameter, (3)~short-range interactions, and
(4)~no additional symmetries belongs to the DP universality class.

% ----------------------------------------------------------------
\subsection{The Three Ingredients}

The simplest lattice model in the DP universality class is the
\emph{contact process} \cite{harris1974}: sites on a grid are active
or inactive; active sites activate neighbors at rate $\lambda$ and
deactivate at rate~1.  The all-inactive state is absorbing.  At
critical $\lambda_c$, the system undergoes a continuous phase
transition with DP exponents.

What are the minimal ingredients for this transition?  Following
Deneubourg et al.~\cite{deneubourg1990} and the ant trail literature,
we identify three:

\begin{enumerate}
    \item \textbf{Amplification} (positive feedback): activity at a
          site increases the probability of future activity at that
          site or its neighbors.
    \item \textbf{Injection} (content-dependent deposition): active
          sites deposit a trace into a shared medium---a signal that
          persists beyond the lifetime of the activity that created it.
    \item \textbf{Erosion} (uniform decay): traces fade at a constant
          rate, preventing runaway accumulation and enabling the system
          to forget.
\end{enumerate}

These three ingredients are necessary and sufficient for the DP phase
transition.  Remove any one and the transition vanishes: without
amplification, traces have no effect; without injection, there is
nothing to amplify; without erosion, the system saturates.

A crucial additional constraint, revealed by systematic study of
our lattice ant model, is that \textbf{erosion must be uncoupled from
the field it erodes}.  Of nine model variants tested, only those with
uniform, field-independent decay produced edge-of-chaos dynamics.
Every variant that introduced decay or noise \emph{correlated} with
the pheromone concentration---follow-probability modulation,
habituation, sensory saturation, coupled deposition---destroyed the
critical ridge entirely.  The mechanism is intuitive: field-correlated
perturbations feed back into the amplification loop, creating a
self-cancelling cycle that prevents the system from reaching the
delicate balance between growth and decay that defines criticality.
In stigmergic systems, the hard problem is not \emph{producing}
edge-of-chaos dynamics but \emph{not destroying} them.  We therefore
treat \textbf{field-independent erosion} as a fourth design
requirement, equal in importance to the three ingredients themselves.

% ----------------------------------------------------------------
\subsection{Ant Trails as the Simplest Illustration}

Ant trail formation is the most intuitive system exhibiting these
ingredients.  Ants deposit pheromone (injection), subsequent ants
preferentially follow stronger trails (amplification), and pheromone
evaporates (erosion).  Deneubourg's bifurcation model
\cite{deneubourg1990} shows that above a critical interaction
frequency, symmetry breaks and coherent trails emerge---a continuous
phase transition with the same mathematical structure as neuronal
avalanches.

We emphasize: the ant model is not the \emph{motivation} for our
architecture.  It is the simplest system that makes the critical
ingredients visible.  The motivation comes from the shared
universality class connecting ant dynamics and neural dynamics.

% ----------------------------------------------------------------
\subsection{Neural Biological Parallels}

Biological neurons possess mechanisms that are functionally
stigmergic---activity leaves traces in a shared medium that modulate
future activity:

\begin{itemize}
    \item \textbf{Synaptic tagging} \cite{frey1997}: A weak stimulus
          creates a short-lived, protein-synthesis-independent ``tag''
          at a synapse.  A strong stimulus at \emph{any} tagged synapse
          triggers plasticity-related protein synthesis captured by
          \emph{all} tagged synapses.  Tags decay without reinforcement
          (erosion), activity creates them (injection), and they bias
          future plasticity (amplification).

    \item \textbf{Nitric oxide volume transmission}
          \cite{bhatt2000,philippides2005}: NO diffuses isotropically
          through tissue, affecting neurons with no active synaptic
          connection to the source.  This is four-dimensional
          environmental signaling---a diffuse field, not a directed
          synapse.

    \item \textbf{Perineuronal nets}: Extracellular matrix structures
          surrounding neurons are modified by experience, physically
          constraining future plasticity---a literal environmental
          trace.

    \item \textbf{Activity-dependent myelination}: Heavily used axons
          receive more myelin, increasing conduction speed.  Use
          modifies the physical substrate, which modifies future signal
          propagation---trail reinforcement.
\end{itemize}

None of these mechanisms are typically called ``stigmergic'' in the
neuroscience literature, but they satisfy the definition precisely:
indirect coordination through modification of a shared environment.
The term \emph{stigmergy} (Grass\'{e}, 1959) was coined for termite
construction, but the mathematical structure---injection, amplification,
erosion leading to a DP-class phase transition---appears to be a
universal motif in biological information processing.

% ----------------------------------------------------------------
\subsection{From Criticality to Architecture}

Our hypothesis: importing the three critical ingredients---plus the
uncoupling constraint---into transformer attention should produce
measurable computational benefits, particularly for tasks requiring
sequential state maintenance---where standard attention's memoryless
query/key matching is provably limited.

In a companion paper \cite{stigmergic2024}, we tested this with the
simplest possible implementation: a \emph{scalar} pheromone field
biasing attention logits.  The results established that the feedback
loop matters (deposits must come from attention \emph{output}, not
input, to close the amplification loop), but the scalar field proved
too weak.  The effect was directional: $+0.07$ accuracy on parity
length generalization, modest improvement on flipflop, no effect on
cumulative sum.  A scalar field decoupled from the attention geometry
cannot express state-dependent routing---it can only say ``attend more
here,'' not ``attend \emph{differently} here.''  Softmax
normalization absorbs most of the scalar signal.

This paper proposes \emph{dual-channel attention}: a vector social
field that modulates the \emph{geometry} of key space, not just the
magnitude of attention scores.  The field operates in the same
representational space as Q/K matching, making it a co-equal routing
mechanism rather than a decorative bias.  Crucially, the social
field's decay is \emph{uniform}---a fixed retention factor $\retain$
applied identically regardless of the field's current magnitude---in
accordance with the uncoupling constraint.  Content-dependent gating
controls what is \emph{deposited}; erosion remains content-blind.


% ================================================================
\section{Why Scalar Biases Are Insufficient}
\label{sec:motivation}

Consider two positions $j_1, j_2$ with raw attention scores $s_1,
s_2$ and scalar field values $\phi_1, \phi_2$.  The relative
attention ratio is:
\begin{equation}
    \frac{a_{j_1}}{a_{j_2}}
    = \frac{\exp(s_1 + \phi_1)}{\exp(s_2 + \phi_2)}
    = \frac{\exp(s_1)}{\exp(s_2)} \cdot \exp(\phi_1 - \phi_2)
\end{equation}
Only the \emph{difference} $\phi_1 - \phi_2$ matters.  If the field
encodes a running state (e.g., parity), all positions receive similar
accumulated values, making the difference small and the bias
negligible.

A scalar bias per position can only say ``attend more here'' or
``attend less here.''  It cannot say:
\begin{itemize}
    \item ``The sequence so far contains an odd count---match against
          keys that encode parity-sensitivity.''
    \item ``The last state-changing event was a \texttt{SET 5}
          command---adjust the key space to surface positions carrying
          that register value.''
\end{itemize}

These require modulating the \emph{geometry} of attention, not its
\emph{magnitude}.  A vector field that modulates keys changes which
\emph{dimensions} of the key space matter, creating orthogonal
routing channels that softmax normalization cannot absorb.


% ================================================================
\section{Why Start with Attention}
\label{sec:why_attention}

The three critical ingredients could in principle be added to any
differentiable computation: MLPs, convolutions, SSMs.  We target
attention specifically for three reasons.

First, attention is \emph{memoryless between positions}.  Each
position computes its output from scratch via Q/K matching, with no
carried state.  This is precisely the property that makes standard
attention fail at sequential state-tracking tasks (parity, counting,
register maintenance)---tasks where biological neural systems, which
\emph{do} have environmental traces (synaptic tags, NO fields),
succeed.  Attention is where the gap between biological and artificial
computation is largest.

Second, attention already has the \emph{amplification} ingredient in
nascent form.  The softmax nonlinearity creates winner-take-all
dynamics: small score differences produce large attention
differences.  What is missing is \emph{injection} (activity leaving a
persistent trace) and \emph{erosion} (traces decaying).  Adding these
two ingredients to a system that already amplifies should be the
minimal intervention needed to reach the critical regime.

Third, attention has a natural \emph{geometric structure}---the Q/K
dot-product space---that a social field can modulate.  Unlike scalar
activations in MLPs, attention keys are vectors in a learned
representation space.  A vector social field can reshape this space,
creating content-and-history-dependent routing that operates in the
same representational language as the structured lookup.  This is
qualitatively different from adding a recurrence alongside attention
(as in Mamba or RWKV): the field is \emph{inside} the attention
computation, not parallel to it.


% ================================================================
\section{Three Architectures}
\label{sec:architecture}

Standard attention is purely \emph{individual}: each position
performs its own content-based lookup, with no memory of what prior
positions found useful.  In the ant analogy, this is an ant following
only its own sensory gradient---functional, but ignorant of the
colony's experience.

We augment attention with a \emph{social field}: a shared,
decaying, vector-valued signal that accumulates deposits from
attention outputs and feeds back into future attention computations.
The social field is the colony's pheromone---an environmental medium
that carries the collective processing history of the sequence.  Each
position reads the social field (what has worked before?) and writes
to it (what did I find?), creating the feedback loop required for
criticality.

We propose three architectures that implement this social field
within attention, ordered by expressivity.  All three share the same
field recurrence but differ in how the field couples to the attention
geometry.  Together they form a progression from the minimal
intervention to a fully content-dependent modulation of the key
manifold.

% ----------------------------------------------------------------
\subsection{Shared Component: The Social Field}

\begin{definition}[Social field]
\label{def:trail_field}
For head $h$, the social field at position $i$ is a vector
$\field_i^{(h)} \in \R^{d_f}$ defined by the recurrence:
\begin{equation}
    \field_i^{(h)} = \retain \cdot \field_{i-1}^{(h)} + \bm{d}_i^{(h)},
    \quad \field_{-1}^{(h)} = \bm{0}
    \label{eq:trail_recurrence}
\end{equation}
where:
\begin{itemize}
    \item $\retain = 1 - \evap \in (0, 1)$ is the retention rate
          (erosion---the field forgets without reinforcement)
    \item $\bm{d}_i^{(h)} \in \R^{d_f}$ is the deposit at
          position $i$, computed from the attention output
          (injection---processing leaves a trace):
          \begin{equation}
              \bm{d}_i^{(h)} = W_{\text{dep}}^{(h)} \bm{o}_i
              \label{eq:vector_deposit}
          \end{equation}
          with $W_{\text{dep}}^{(h)} \in \R^{d_f \times d}$ and
          $\bm{o}_i$ the attention output at position $i$
    \item $d_f$ is the \emph{field dimension}, controlling the
          information capacity of the social channel
\end{itemize}
The field implements two of the three critical ingredients directly.
Amplification---the third ingredient---arises from the coupling: the
field modulates attention, which changes the output, which changes
the deposit, which changes the field.
\end{definition}

The term \emph{social} is deliberate.  In attention, queries and keys
are \emph{individual}---each position's private representation.  The
social field is \emph{collective}---it belongs to no single position
but accumulates the processing history of the entire sequence so far.
An individual position reads the social field to discover what prior
processing found useful and writes to it to inform future positions.
This is precisely the stigmergic coordination loop: indirect
communication through a shared, decaying environmental medium.

The field dimension $d_f$ controls state capacity.  In the scalar
case ($d_f = 1$), this reduces to the stigmergic attention bias of
our companion paper.  We set $d_f = d_h$ (the head dimension) so the
social field and keys live in the same representational space.

\begin{remark}[State capacity]
Each head carries $d_f$ continuous values of social state.  With $H$
heads, the total social state is $H \times d_f$---compared to 0 for
standard attention and $n$ for SSMs with state dimension $n$.  For
$H = 4, d_f = 12$, this is 48 values of persistent state, vastly
more than the $H = 4$ scalars of the stigmergic architecture.
\end{remark}

\begin{remark}[Why modulate keys, not queries?]
\label{rem:keys_not_queries}
Keys are the \emph{offerings} of past positions; queries are the
\emph{demands} of the current position.  The social field modulates
offerings, not demands: it reshapes \emph{what information is
available} based on collective processing history.  This preserves
the individual nature of queries---each position's sensory state
remains purely content-driven---while making the key landscape
socially informed.  In the ant analogy: pheromone marks where ants
\emph{have walked} (keys), not where they \emph{want to go}
(queries).
\end{remark}

% ----------------------------------------------------------------
\subsection{Architecture I: Additive Key Shift}
\label{sec:arch_additive}

The simplest proposal.  The social field is projected to key dimension
and added to keys before Q/K matching:
\begin{equation}
    \K_j'^{(h)} = \K_j^{(h)} + W_{\text{mod}}^{(h)} \field_j^{(h)}
    \label{eq:additive_mod}
\end{equation}
where $W_{\text{mod}}^{(h)} \in \R^{d_h \times d_f}$.

\paragraph{Interpretation.}  The attention score decomposes cleanly
into two channels:
\begin{align}
    \text{score}(i, j) &= \frac{1}{\sqrt{d_h}} \q_i^{(h)\top}
    \K_j'^{(h)} \notag \\
    &= \underbrace{\frac{\q_i^{(h)\top} \K_j^{(h)}}{\sqrt{d_h}}}_%
    {\text{structured (content match)}}
    + \underbrace{\frac{\q_i^{(h)\top} W_{\text{mod}}
    \field_j^{(h)}}{\sqrt{d_h}}}_{\text{social (history match)}}
    \label{eq:score_decomp}
\end{align}
Unlike the scalar bias of stigmergic attention, the social term is a
\emph{full dot product} between the query and a field-derived
vector---it cannot be absorbed by softmax normalization because it
varies across the $d_h$ dimensions.

\paragraph{Parameters.}  $W_{\text{dep}}$ ($d_f \times d$ per head)
+ $W_{\text{mod}}$ ($d_h \times d_f$ per head).  With $d_f = d_h$:
$\sim\!25\%$ overhead on the attention layer.

\paragraph{Why start here.}  This is the minimal intervention that
gives the field access to the geometric structure of attention.  If
additive modulation shows no effect, the problem is not expressivity
but something more fundamental (e.g., the sequential bottleneck
prevents learning, or the tasks don't benefit from social state).

% ----------------------------------------------------------------
\subsection{Architecture II: Gating Modulation}
\label{sec:arch_gating}

The social field produces per-dimension gates that amplify or suppress
key dimensions:
\begin{equation}
    \K_j'^{(h)} = \K_j^{(h)} \odot \sigma\!\bigl(
        W_{\text{gate}}^{(h)} \field_j^{(h)} + \bm{b}_g\bigr)
    \label{eq:gating_mod}
\end{equation}
where $\sigma$ is the sigmoid function and $\odot$ is elementwise
multiplication.

\paragraph{Interpretation.}  Rather than shifting keys, gating
\emph{selects} which dimensions of the key space are active.  The
field controls which semantic features are matchable at each
position.  A gate near 0 on dimension $k$ effectively removes that
feature from the key; a gate near 1 preserves it.  This is a
multiplicative interaction---qualitatively stronger than additive
because it can \emph{silence} dimensions, not just nudge them.

\paragraph{Parameters.}  Same as additive plus a bias vector
$\bm{b}_g \in \R^{d_h}$ per head.  Negligible overhead beyond
Architecture~I.

\paragraph{When to prefer over I.}  Gating is appropriate when the
task requires \emph{selective attention}---attending to different
aspects of keys depending on accumulated state.  For example: in
bracket matching, the field might gate out ``nesting depth''
dimensions when the stack is empty and gate them in when brackets
are open.

% ----------------------------------------------------------------
\subsection{Architecture III: Rotational Modulation}
\label{sec:arch_rotational}

The field produces a content-and-state-dependent rotation of the key
vector:
\begin{equation}
    \K_j'^{(h)} = \K_j^{(h)} + W_{\text{rot}}^{(h)}
    (\field_j^{(h)} \otimes \K_j^{(h)})
    \label{eq:rotational_mod}
\end{equation}
where $\otimes$ denotes a bilinear interaction (outer product
projected back to $\R^{d_h}$) and
$W_{\text{rot}}^{(h)} \in \R^{d_h \times d_f \cdot d_h}$.

\paragraph{Interpretation.}  This is the most expressive variant.
The modulation depends on \emph{both} the field state and the key
content---different keys at the same position are rotated differently.
This creates a content-and-history-dependent routing that can
implement arbitrary learned transformations of the key manifold.

\paragraph{Parameters.}  $W_{\text{rot}}$ is $d_h \times (d_f \cdot
d_h)$ per head, which is $O(d_h^2 \cdot d_f)$.  This is expensive:
for $d_h = d_f = 12$, it is $12 \times 144 = 1728$ parameters per
head versus 144 for additive.  A low-rank factorization
$W_{\text{rot}} = U V^\top$ with rank $r \ll d_f$ can reduce this.

\paragraph{When to prefer over I--II.}  Rotational modulation is
appropriate when the task requires the field to implement different
transformations depending on key content---e.g., when the same
accumulated state should affect ``parity-encoding'' keys differently
from ``register-encoding'' keys.  We expect this to matter primarily
at scale, where different heads specialize.

% ----------------------------------------------------------------
\subsection{Shared Forward Pass}

All three architectures share the same forward pass structure.  At
each position $i$:
\begin{enumerate}
    \item \textbf{Project} input to queries, keys, values:
          $\q_i, \K_i, \vv_i = \text{QKV}(\x_i)$
    \item \textbf{Read social field}: retrieve accumulated
          $\field_j^{(h)}$ for all $j < i$
    \item \textbf{Modulate keys}: $\K_j'^{(h)} = f(\K_j^{(h)},
          \field_j^{(h)})$ using one of the three strategies
    \item \textbf{Attend}: compute attention weights using modulated
          keys
          \begin{equation}
              a_{ij}^{(h)} = \text{softmax}_j\!\Bigl(
              \frac{\q_i^{(h)\top} \K_j'^{(h)}}{\sqrt{d_h}}
              \Bigr) \quad j \leq i
          \end{equation}
    \item \textbf{Aggregate}: $\bm{o}_i^{(h)} = \sum_{j \leq i}
          a_{ij}^{(h)} \vv_j^{(h)}$
    \item \textbf{Deposit}: $\bm{d}_i^{(h)} = W_{\text{dep}}^{(h)}
          \bm{o}_i$ \quad (injection)
    \item \textbf{Update field}: $\field_i^{(h)} = \retain \cdot
          \field_{i-1}^{(h)} + \bm{d}_i^{(h)}$ \quad (erosion)
\end{enumerate}

The feedback loop closes the amplification ingredient:
\begin{equation}
    \field \xrightarrow{\text{modulate } K}
    \text{attention} \xrightarrow{\text{aggregate}}
    \text{output} \xrightarrow{W_{\text{dep}}}
    \text{deposit} \xrightarrow{\text{accumulate}} \field
\end{equation}

\begin{table}[h]
\centering
\small
\begin{tabular}{@{}lccc@{}}
\toprule
& \textbf{I: Additive} & \textbf{II: Gating}
& \textbf{III: Rotational} \\
\midrule
Modulation type & shift & selection & rotation \\
Field--key interaction & linear & multiplicative & bilinear \\
Extra params / head & $d_h d_f + d_f d$ & same + $d_h$ &
$d_h^2 d_f + d_f d$ \\
Content-dependent & no & no & yes \\
Can silence dimensions & no & yes & yes \\
Expressivity & low & medium & high \\
\bottomrule
\end{tabular}
\caption{Summary of the three proposed architectures.  All share the
same social field recurrence and forward pass; they differ only in
how the field modulates keys.  We recommend starting experiments
with Architecture~I (additive) as the simplest test of whether
social fields provide computational benefit.}
\label{tab:architectures}
\end{table}

% ----------------------------------------------------------------
\subsection{The Individual/Social Decomposition}

Standard attention performs a single type of lookup---purely
individual, purely content-based:
\begin{equation}
    \text{score}(i, j) = \q_i^\top \K_j / \sqrt{d_h}
    \quad \text{(``what content at $j$ matches my query?'')}
\end{equation}

All three architectures decompose this into an individual channel
and a social channel.  For Architecture~I (additive), the
decomposition is exact (cf.~\cref{eq:score_decomp}):
\begin{align}
    \text{individual}(i, j) &= \q_i^\top \K_j / \sqrt{d_h}
    \quad \text{(``what content matches?'')}
    \label{eq:structured} \\
    \text{social}(i, j) &= \q_i^\top
    (W_{\text{mod}} \field_j) / \sqrt{d_h}
    \quad \text{(``what has the sequence learned?'')}
    \label{eq:social}
\end{align}
The individual channel is memoryless: it depends only on the current
query and the static key.  The social channel carries history: it
depends on the accumulated social field, which encodes what prior
attention computations found productive.  For Architectures~II and
III, the decomposition is implicit but the principle is the same:
one channel routes by content, the other by collective experience.

This is the core architectural claim: attention conflates individual
and social information processing into a single dot product.
Separating them---giving each its own mechanism, its own
parameters, its own representational capacity---should allow the
model to learn both more effectively.


% ================================================================
\section{Relationship to Existing Architectures}
\label{sec:related}

\begin{table}[h]
\centering
\small
\begin{tabular}{@{}lcccc@{}}
\toprule
\textbf{Architecture} & \textbf{Modulation} & \textbf{State}
& \textbf{Feedback} & \textbf{Mechanism} \\
\midrule
Standard Transformer & none & --- & --- & Q/K matching \\
ALiBi \cite{press2022alibi} & logit bias & --- & --- & fixed distance penalty \\
T5 relative bias & logit bias & --- & --- & learned per-distance \\
RoPE \cite{su2024rope} & query/key rotation & --- & --- & fixed rotation \\
\addlinespace
Stigmergic (scalar) & logit bias & $H$ scalars & nonlinear & output $\to$ scalar field \\
\addlinespace
RetNet \cite{sun2023retnet} & key decay & $d$ & linear & $\gamma$-mask on keys \\
Mamba \cite{gu2023mamba} & separate path & $n$ & linear & selective SSM \\
RWKV \cite{peng2023rwkv} & separate path & $d$ & linear & linear attention \\
Feedback Transformer & extra KV & $d$ & nonlinear & top-layer output \\
\addlinespace
\textbf{Dual-channel} & \textbf{key modulation} & $\bm{H \times d_f}$
& \textbf{nonlinear} & \textbf{output $\to$ vector field} \\
\bottomrule
\end{tabular}
\caption{Architectural comparison.  Dual-channel attention is unique
in using a \emph{vector} recurrence to modulate \emph{keys}, creating
a content-dependent routing channel that operates in the same
representational space as Q/K matching.}
\label{tab:comparison}
\end{table}

\paragraph{ALiBi and RoPE} modify attention geometry but
\emph{statically}---the modification depends on position, not on
content or processing history.  Dual-channel attention is ``ALiBi
where the bias is a learned vector field driven by what attention
has found useful.''

\paragraph{RetNet} applies a multiplicative $\gamma^{i-j}$ decay to
keys, which is structurally similar to our decaying field.  But
RetNet's decay is \emph{fixed and content-independent}---it provides
exponential forgetting but no feedback.  The social field's deposits
are content-dependent and output-derived, creating adaptive decay.

\paragraph{Mamba and RWKV} maintain recurrent state but in a
\emph{separate pathway} from attention.  The SSM and attention do not
interact---they are stacked, not coupled.  In dual-channel attention,
the recurrent state (social field) directly modulates the attention
computation, and the attention output directly feeds the recurrence.
The two mechanisms are co-dependent, not independent.

\paragraph{Feedback Transformer \cite{fan2021feedback}} is the
closest prior work.  It feeds the top layer's output as an
additional key-value pair for the next position.  This provides
output feedback but through a \emph{single vector} added to the
key-value set, not through modulation of \emph{all} keys.  The
social field modulates every key at every position through a
decaying recurrence, providing graded, history-dependent routing
rather than a single feedback signal.

\paragraph{What is novel.}
The specific combination---a vector recurrence that modulates keys
via additive/gating/rotational transforms, with deposits derived
from attention output---does not appear in prior work to our
knowledge.  More fundamentally, the architecture is not motivated by
analogy but by universality: the three ingredients
(amplification, injection, erosion) are the minimal recipe for DP-class
phase transitions, and we import them into attention to test whether
criticality confers computational advantages in learned sequence
processing, as it does in biological neural systems
\cite{beggs2003,shew2013}.


% ================================================================
\section{Analysis}
\label{sec:analysis}

% ----------------------------------------------------------------
\subsection{Parameter Cost}

For additive modulation with $d_f = d_h$:
\begin{itemize}
    \item Deposit projection: $W_{\text{dep}} \in \R^{d_f \times d}$,
          adding $d_f \times d$ parameters per head
    \item Modulation projection: $W_{\text{mod}} \in
          \R^{d_h \times d_f}$, adding $d_h \times d_f$ per head
    \item Total per head: $d_f(d + d_h)$. With $d_f = d_h = d/H$:
          $\frac{d}{H}\bigl(d + \frac{d}{H}\bigr)
          = \frac{d^2}{H} + \frac{d^2}{H^2}$
    \item Total for $H$ heads: $d^2 + d^2/H$
    \item Compared to QKV projection ($3d^2$) and output projection
          ($d^2$): the social field adds $\sim 25\%$ parameters to
          the attention layer (for typical $H$)
\end{itemize}

Gating adds bias terms ($d_h$ per head = $d$ total).  Rotational
adds a bilinear interaction, roughly $d_h^2 \times d_f$ per head,
which can be expensive.

% ----------------------------------------------------------------
\subsection{Computational Cost}

The sequential bottleneck remains: the feedback loop requires
position-by-position processing because position $i$'s deposit
depends on the field state from positions $0, \ldots, i-1$.

\paragraph{Training.}  During training on sequences of length $T$,
the forward pass requires a sequential scan of $T$ steps.  At each
step: one matrix-vector multiply for the deposit ($O(d_f \times d)$),
one for the modulation ($O(d_h \times d_f)$), and a vector add for
the field update ($O(d_f)$).  The attention computation itself
($O(T \times d_h)$ per position) dominates.  Total cost is
$O(T^2 d_h)$ same as standard attention, but the constant factor
increases and parallelism across positions is lost.

\paragraph{Inference.}  In autoregressive inference, standard
attention is already sequential.  The social field adds only
$O(d_f)$ work per position for the field update and $O(d_h
\times d_f)$ for key modulation---negligible compared to the
$O(T \times d_h)$ attention computation.  \textbf{At inference
time, the social field is essentially free.}

\paragraph{Potential acceleration.}  The sequential constraint
during training could be relaxed via:
\begin{itemize}
    \item \emph{Chunked processing}: process chunks of $C$ positions
          in parallel (using only the field state at the chunk
          boundary), then update the field.  This trades accuracy for
          parallelism.
    \item \emph{Truncated backpropagation}: detach the field from the
          computation graph every $K$ positions, allowing parallel
          chunks of length $K$ with approximate gradients.
    \item \emph{Linear attention within chunks}: use linear attention
          (which supports parallel scan) for the field-dependent
          component, keeping softmax attention for the structured
          channel.
\end{itemize}

% ----------------------------------------------------------------
\subsection{Length Generalization}

The social field recurrence $\field_i = \retain \cdot \field_{i-1}
+ \bm{d}_i$ is position-independent: it depends only on the
previous field state and the current deposit, not on absolute
position.  This means the field dynamics are \emph{identical} at
position 50 during training and position 500 during inference---the
recurrence does not know where it is in the sequence.

Combined with a position-independent attention mechanism (RoPE
or no positional encoding), the entire dual-channel system is
position-independent.  The structured channel matches by content,
the social channel routes by accumulated history, and neither
requires knowing the absolute position.  This is the architectural
basis for length generalization.

Our prior work showed that even a scalar field (stigmergic
attention) provides measurable length generalization on parity,
despite using learned positional embeddings that break at unseen
positions.  With vector key modulation and position-independent
encoding, the advantage should be substantially larger: the field
carries more state (vectors vs scalars), modulates geometry (not
just magnitude), and the attention mechanism itself still functions
at OOD lengths.


% ================================================================
\section{Experimental Design}
\label{sec:experiments}

We propose experiments targeting the specific capabilities that
dual-channel attention provides beyond standard attention:
sequential state tracking and length generalization.

\subsection{Tasks}

\paragraph{Parity (1-bit state).}
The simplest state-tracking task, previously shown to benefit from
scalar feedback.  Dual-channel attention should show a larger effect
because the vector field can encode parity as a \emph{direction} in
key space, not just a scalar magnitude.

\paragraph{Cumulative sum mod $V$ ($\log_2 V$-bit state).}
Previously failed with scalar field because the state exceeds scalar
capacity.  With $d_f = d_h = 12$ and $H = 4$ heads, the field has
48 continuous values of state---far more than the 4 bits needed for
mod-16 arithmetic.

\paragraph{Bracket matching (counter + stack).}
Given sequences of open/close brackets with possible nesting, predict
whether the sequence is balanced at each position.  This requires a
counter (nesting depth) and implicit stack (which bracket type to
match).  The field could encode depth as a vector magnitude and
bracket type as a vector direction.

\paragraph{Evaluation protocol.}
For all tasks: train on sequences of length $T$, evaluate on
$T, 2T, 3T, 4T$.  Use RoPE or no positional encoding to isolate
the field's contribution to length generalization.  Compare:
standard attention, scalar stigmergic attention, and dual-channel
attention with additive and gating modulation.


% ================================================================
\section{Scaling Considerations}
\label{sec:scaling}

\paragraph{The sequential bottleneck is the primary concern.}
At training time, the feedback loop forces sequential processing.
For long sequences ($T > 1000$), this makes training prohibitively
slow compared to standard parallel attention.  However:
\begin{itemize}
    \item At inference time (autoregressive generation), processing
          is already sequential---the social field adds negligible cost.
    \item For encoder models processing fixed-length inputs, the
          sequential constraint applies, but can be mitigated by
          chunked processing.
    \item The field could be applied to only a subset of layers
          (e.g., every 4th layer), reducing the sequential burden
          while maintaining most of the benefit.
\end{itemize}

\paragraph{Field dimension scaling.}
Setting $d_f = d_h$ means the field scales with model size
automatically.  A model with $d = 768, H = 12$ gets $d_f = 64$
per head, providing $12 \times 64 = 768$ values of persistent
state.  This is comparable to Mamba's state dimension and
should be sufficient for complex state tracking.

\paragraph{Gradient flow through the recurrence.}
The field recurrence $\field_i = \retain \field_{i-1} + \bm{d}_i$
has gradient $\partial \field_i / \partial \field_j = \retain^{i-j}$.
For $\retain = 0.95$ and $i - j = 100$: $0.95^{100} \approx 0.006$.
Gradients through the field decay exponentially, potentially causing
vanishing gradient problems for long-range dependencies.  Possible
mitigations:
\begin{itemize}
    \item \emph{Learnable retention}: let $\retain$ be a learned
          parameter per head, initialized near 1.0
    \item \emph{Multi-scale fields}: different heads at different
          retention rates, spanning timescales
    \item \emph{Gated update}: $\field_i = g_i \odot \field_{i-1}
          + (1 - g_i) \odot \bm{d}_i$ with learned gate $g_i$,
          allowing selective persistence (cf.\ GRU)
\end{itemize}

\paragraph{Interaction with existing optimizations.}
Flash Attention exploits the fact that softmax attention is
block-decomposable.  Key modulation preserves this property---the
modulated keys $\K'$ can be precomputed (for non-feedback variants)
or computed per-chunk (for chunked feedback).  KV caching at
inference time extends naturally: cache the \emph{modulated} keys
$\K'_j$ along with values.


% ================================================================
\section{Discussion}
\label{sec:discussion}

\subsection{The Individual/Social Decomposition}

The core claim is that attention performs two conceptually distinct
operations that benefit from separate mechanisms:

\begin{enumerate}
    \item \textbf{Individual lookup}: ``Given what I'm looking for
          (query) and what's available (key), how relevant is each
          position?''  This is \emph{content-addressed} and
          \emph{memoryless}---it depends only on the current
          position's private state.

    \item \textbf{Social routing}: ``Given the collective processing
          history of the sequence, which directions in key space are
          worth attending to?''  This is \emph{history-dependent}
          and \emph{collectively derived}---it comes from the social
          field, not from any single position.
\end{enumerate}

Standard attention collapses these into a single Q/K dot product.
Learned positional encodings partially separate them (position
provides structure-dependent routing) but are static and
length-limited.  The social field provides \emph{dynamic},
\emph{experience-derived}, \emph{length-independent} routing that
adapts to each sequence individually---the collective intelligence
of the sequence's own processing history.

\subsection{What the Social Field Learns}

In biological stigmergic systems, environmental traces encode three
types of information:

\begin{enumerate}
    \item \textbf{Recency}: recent deposits are stronger (exponential
          decay).  In the social field: $\retain^{i-j}$ weighting.
    \item \textbf{Frequency}: frequently productive paths accumulate
          more signal.  In the social field: positions that
          consistently produce useful outputs accumulate deposits.
    \item \textbf{Direction}: field gradients indicate which way to
          go.  In the vector social field: the field encodes a
          \emph{direction} in key space, not just a magnitude.
\end{enumerate}

The scalar stigmergic field of our companion paper captured recency
and frequency but not direction.  The vector social field captures
all three, which is why we expect qualitatively different behavior:
the social channel can steer attention through a learned geometry,
not merely amplify or suppress it.

\subsection{Limitations and Open Questions}

\begin{itemize}
    \item \textbf{Training speed}: The sequential bottleneck is real
          and significant.  Chunked approximations trade accuracy for
          speed, and the optimal tradeoff is unknown.
    \item \textbf{Task generality}: State-tracking tasks are
          constructed to favor recurrent mechanisms.  Whether the
          social field helps on natural language tasks---where the
          ``state'' is diffuse and high-dimensional---is unknown.
    \item \textbf{Optimization difficulty}: The feedback loop creates
          long gradient paths through the recurrence and the
          attention mechanism.  This may make training unstable or
          slow to converge.  The seed-dependent bimodality observed
          in scalar experiments suggests the optimization landscape
          has multiple basins.
    \item \textbf{Redundancy with existing mechanisms}: Multi-layer
          transformers already build up sequential context through
          layer composition.  The social field may provide redundant
          state-tracking that deeper networks achieve anyway.
\end{itemize}

\subsection{Isn't the Residual Stream Already a Social Field?}
\label{sec:already_social}

A natural objection: deep transformers already exhibit the three
critical ingredients, so why add another mechanism?  The residual
stream \emph{is} a shared medium.  Each layer deposits into it
(injection).  Softmax creates winner-take-all dynamics that compound
across layers (amplification).  LayerNorm rescales activations,
preventing runaway accumulation (erosion of a sort).  The critic
concludes: the transformer is already a stigmergic system, and the
social field reinvents what depth provides for free.

This objection is half right, and engaging it precisely clarifies what
the social field actually contributes.

\paragraph{The residual stream is social across layers, not across
positions.}
At any given layer, each position's attention computation is memoryless
with respect to other positions' \emph{processing outcomes}.
Position~$i$ sees the raw keys and values of positions $j < i$, but it
does not see what positions $j\!-\!1, j\!-\!2, \ldots$ found useful
when \emph{they} attended.  The residual stream carries information
forward through depth (layers) but not through sequence (positions
within a layer).  The social field carries information forward through
the sequence dimension---precisely the dimension where standard
attention has no memory.  This is why state-tracking tasks (parity,
counting, register maintenance) are hard for standard transformers:
they must emulate sequential accumulation using combinatorial
depth$\,\times\,$width, rather than having an architectural primitive
for it.

\paragraph{Depth is fixed; the social field scales with sequence
length.}
A 12-layer transformer has exactly 12 opportunities for inter-layer
coordination, whether the input is 10 tokens or 10{,}000.  The social
field accumulates deposits proportional to sequence length---it has
\emph{more} collective memory for longer sequences.  This is the
architectural basis for length generalisation: the field's recurrence
is position-independent and naturally extends to unseen lengths, while
depth-based state emulation is capacity-bounded.

\paragraph{The residual stream does not modulate the geometry of
attention.}
Information in the residual stream enters the next layer's attention
through generic QKV projections.  The stream changes what gets
projected, but through a learned linear transform that treats the
residual vector as undifferentiated input.  The social field
specifically modulates \emph{keys}---it operates in the same
representational space as the Q/K dot product, reshaping the matching
geometry rather than merely providing different input features.  This
is the difference between ``information flows forward'' (trivially true
of any deep network) and ``a decaying field reshapes which directions
in key space are matchable.''

\paragraph{The uncoupling constraint: where the ant model builds
intuition.}
The most non-obvious distinction concerns erosion.  LayerNorm's
rescaling is \emph{field-correlated}: it responds to the magnitude of
the activations it normalises, shrinking large values and expanding
small ones.  Our lattice ant study provides direct intuition for why
this matters.  Of nine model variants we tested, every variant whose
decay was correlated with the pheromone field---follow-probability
modulation, habituation, sensory saturation---destroyed the critical
ridge entirely.  Only variants with uniform, field-independent decay
preserved edge-of-chaos dynamics.  The mechanism is intuitive once seen
in the ant system: field-correlated perturbations feed back into the
amplification loop, creating a self-cancelling cycle that prevents the
delicate balance between growth and decay that defines criticality.
LayerNorm does exactly this---it is a field-correlated erosion that
actively counteracts the very accumulation patterns that would
otherwise self-organise.  The social field's uniform $\retain$ decay,
by contrast, satisfies the uncoupling constraint.  This distinction
would be purely theoretical without the ant model; the lattice
simulations make it concrete and testable.

\paragraph{The honest assessment.}
Multi-layer transformers \emph{do} track state through depth---modern
large models handle parity, counting, and bracket matching.  But they
do so by learning to emulate sequential state maintenance, not by
having a native mechanism for it.  The social field makes sequential
state explicit.  Whether this leads to faster learning, better
generalisation, or mere redundancy is exactly what the experimental
programme (\cref{sec:experiments}) is designed to test.  The
architectural argument is that three specific properties---positional
memory, uncoupled erosion, and geometric modulation---are absent from
the residual stream and present in the social field.  Whether those
properties \emph{matter} computationally is an empirical question.

\subsection{Lessons from the Ant Model: Uncoupled Fields and Criticality}

A striking result from our lattice ant study is that
\emph{edge-of-chaos dynamics require the decay (erosion) channel to be
uncoupled from the pheromone field itself}.  Of nine model variants
tested, only those with uniform, field-independent decay produced a
critical ridge; every variant that introduced disorder correlated with
the pheromone concentration---follow-probability modulation,
habituation, sensory saturation, coupled deposition---destroyed
criticality by feeding field-dependent noise back into the
amplification loop.  In short: the hard problem is not \emph{producing}
edge-of-chaos dynamics but \emph{not destroying} them with
field-correlated perturbations.

This has a direct architectural implication for dual-channel attention.
The social field's decay parameter $\retain$ should act as a
\emph{uniform} exponential forgetting factor, independent of the
field's current magnitude.  If decay were made content-dependent
(e.g.\ stronger erasure where the field is large, or gated by the
field's own values), the ant model predicts that the system would lose
its ability to self-organise near criticality.  The separation of
concerns---individual attention handles content-dependent computation,
while the social field evolves under content-\emph{independent}
erosion---may therefore be essential, not merely convenient.  Uniform
decay preserves the directed-percolation structure that enables
critical dynamics; field-coupled decay collapses it.


\subsection{Reflections: Sociality as Structure Discovery}

A useful way to distinguish the individual and social channels is by
what kind of knowledge they encode.  Individual lookup is literal: given
this query, find the matching key---a direct content-addressed
retrieval.  The social field does something qualitatively different.  By
accumulating deposits across positions, it builds a \emph{compressed
record of which directions in key space were collectively productive}.
This is not a lookup table but a structural prior: things shaped like
\emph{this} tend to benefit from rules shaped like \emph{that}.  Where
individual attention asks ``what should I do here?'', the social field
asks ``what kind of situation is this, and what has worked for situations
like it?''

This is, in a sense, the hard version of the problem.  Finding literal
associations is cheap; discovering that two superficially different
contexts share hidden structure---that they are instances of the same
\emph{metaphor}---requires the system to compress over its own
processing history.  But the hypothesis motivating this architecture is
that the world is full of such reusable structure, and this is precisely
what makes it learnable.  Biological systems exploit this ruthlessly:
evolution builds on older solutions without relearning from scratch,
reusing regulatory motifs, signalling cascades, and body plans across
wildly different contexts.  The social field is a minimal mechanism for
the same principle---finding and propagating structural regularities
rather than storing each situation individually.

% ================================================================
\section{Conclusion}
\label{sec:conclusion}

We have argued that the three ingredients governing critical phase
transitions in the directed percolation universality
class---amplification, injection, erosion---appear in both ant trail
formation and biological neural plasticity mechanisms (synaptic
tagging, volume transmission, activity-dependent myelination).  The
ant trail model is the simplest system that makes these ingredients
and their phase transition visible; we use it as an illustration,
not a metaphor.

Our companion paper \cite{stigmergic2024} tested the simplest
possible import of these ingredients into attention: a scalar
pheromone field biasing logits.  The results confirmed that the
feedback loop matters (deposits must come from output, closing the
amplification loop) but showed that a scalar field decoupled from
key geometry is too weak---softmax normalization absorbs most of
the signal.  This motivated the present architecture.

Dual-channel attention decomposes the attention mechanism into
individual lookup (content-based Q/K matching) and social routing
(vector field key modulation).  The social field provides
intrinsic sequential memory---$H \times d_f$ continuous values of
persistent, decaying, content-dependent state per layer---without
introducing a separate recurrent pathway.  It operates in the same
representational space as Q/K matching, modulating the \emph{geometry}
of attention rather than its magnitude.

The deeper question is whether criticality plays a functional role
in learned sequence processing---whether there exists a regime where
the social field's feedback dynamics self-organize to a critical
point, and whether this criticality confers the same computational
advantages (maximal dynamic range, information transmission,
sensitivity) observed in biological neural systems at criticality.
The architecture is designed to make this question empirically
testable.


\appendix

\section{Research Notes: Unifying Learning and Inference}
\label{app:learning-inference}

\subsection{The Separation in Standard Attention}

In a vanilla transformer, there is a clean division of labour:
\emph{training} learns the projections $W_Q$, $W_K$, $W_V$---essentially
a filing system for where to put information and how to retrieve it.  At
\emph{inference} time, those projections are frozen, and the model
executes retrieval.  ``In-context learning'' in standard transformers is
somewhat misleadingly named: the model is not learning anything during
the forward pass.  It is recognising patterns it already learned to
recognise during training.  The weights do not change.  It is retrieval
all the way down.

\subsection{What the Social Field Changes}

The social field actually \emph{does} modify itself during the forward
pass.  As it processes positions $1, 2, 3, \ldots$, it accumulates
deposits that reshape the attention geometry for all subsequent
positions.  Position~50 attends in a key space that has been sculpted by
the collective processing outcomes of positions~1--49.  This is
genuinely different: the field is building temporary structure
\emph{during inference} that did not exist before and was not hardcoded
in the weights.

Standard attention learns \emph{where to file}, then retrieves.  The
social field learns \emph{what structure exists in this particular
sequence} as it processes it, and that structure modifies how it
processes what comes next.  The deposit--accumulate--modulate loop is
functionally a learning algorithm operating at inference time.

\subsection{Three Timescales}

The deeper point concerns what ``learning'' means at different
timescales:

\begin{enumerate}
    \item \textbf{Training} (slow): learn $W_Q$, $W_K$, $W_V$,
          $W_{\mathrm{deposit}}$---the rules of the game.
    \item \textbf{Social field accumulation} (fast): discover the
          structure of \emph{this particular sequence}---the shape of
          this game.
    \item \textbf{Individual attention} (instantaneous): make a single
          retrieval given the current state.
\end{enumerate}

Standard transformers have only the slow and instantaneous timescales.
The social field introduces the \emph{middle} timescale---slower than a
single attention operation, faster than weight updates.  This is exactly
the timescale that is missing, and it is the timescale at which
biological plasticity operates: synaptic tagging, short-term
potentiation, neuromodulatory volume transmission---all happening during
``inference'' in the brain.

\subsection{The Biological Parallel}

Biological brains do not have a train/infer separation at all---plasticity
is continuous.  What they \emph{do} have is multiple timescales of
plasticity: fast (ion channel states, milliseconds), medium (synaptic
tagging, minutes to hours), slow (structural plasticity, days to weeks).
The social field is a minimal version of this: it adds one intermediate
plasticity timescale to an architecture that otherwise only has ``frozen
weights'' and ``instantaneous computation.''

The presupposition that the world has countable order is important here.
The social field can only help if the structure it discovers early in a
sequence is \emph{predictive} of structure later in the sequence---which
is exactly the assumption that the world has reusable regularities.  If
every position were independently random, the social field would
accumulate noise.  That it accumulates \emph{signal} is a bet on the
world being structured, and this is the same bet that makes evolution
possible: old solutions remain useful because the world's structure is
stable enough to reuse.

\subsection{A Testable Provocation}

The social field brings learning and inference closer, but they are not
fully unified---the field resets between sequences.  True unification
would mean the field persists and the system improves across sequences,
not just within one.  The social field is more like ``inference that
remembers its own recent history'' than ``learning.''  But this may be
exactly the right intermediate step, and it is notable that biological
short-term plasticity also resets: synaptic facilitation decays, working
memory clears---and yet it is clearly doing something different from pure
retrieval.

The really provocative version of this claim: \emph{the reason
transformers need so many parameters and so much training data is that
they are forced to compile all sequential structure into static weights,
because they lack a mechanism for discovering structure at inference
time}.  The social field could, in principle, let a smaller model with
less training discover structure on-the-fly that a standard transformer
would need to have memorised.  This is testable, and is the direction we
are chipping away at.

\paragraph{Note: multiscale validation on natural language.}
Preliminary experiments training the vector social field on TinyStories
(11.7M parameter GPT, 512-token context, chunk size 128) reveal that the
single-timescale architecture described here---while functional---is
severely limited by its fixed decay rate ($\evap = 0.05$, half-life
$\sim$13 tokens).  Extending the architecture with a second, slow field
(learnable retention converging to half-life $\sim$120 tokens) produces a
36.6\% perplexity increase at late positions when ablated, compared to
2.18\% for the single field, and is the only variant to improve aggregate
perplexity.  The model spontaneously discovers scale separation, with the
slow field spanning the full context window while the fast field retains
local dynamics.  These results, detailed in the companion hierarchical
coupling paper, confirm that the social field's capacity for long-range
influence is real but requires temporal scale structure to be drawn out
by standard next-token training.


\bibliographystyle{plain}
\begin{thebibliography}{99}

% --- Criticality and universality ---

\bibitem{beggs2003}
J.~M.~Beggs and D.~Plenz.
Neuronal avalanches in neocortical circuits.
\emph{Journal of Neuroscience}, 23(35):11167--11177, 2003.

\bibitem{shew2013}
W.~L.~Shew and D.~Plenz.
The functional benefits of criticality in the cortex.
\emph{The Neuroscientist}, 19(1):88--100, 2013.

\bibitem{chialvo2010}
D.~R.~Chialvo.
Emergent complex neural dynamics.
\emph{Nature Physics}, 6:744--750, 2010.

\bibitem{munoz2018}
M.~A.~Mu\~{n}oz.
Colloquium: Criticality and dynamical scaling in living systems.
\emph{Reviews of Modern Physics}, 90:031001, 2018.

\bibitem{janssen1981}
H.-K.~Janssen.
On the nonequilibrium phase transition in reaction-diffusion systems
with an absorbing stationary state.
\emph{Zeitschrift f\"{u}r Physik B}, 42:151--154, 1981.

\bibitem{grassberger1982}
P.~Grassberger.
On phase transitions in Schl\"{o}gl's second model.
\emph{Zeitschrift f\"{u}r Physik B}, 47:365--374, 1982.

\bibitem{harris1974}
T.~E.~Harris.
Contact interactions on a lattice.
\emph{The Annals of Probability}, 2(6):969--988, 1974.

\bibitem{hesse2014}
J.~Hesse and T.~Gross.
Self-organized criticality as a fundamental property of neural systems.
\emph{Frontiers in Systems Neuroscience}, 8:166, 2014.

% --- Ant trails and stigmergy ---

\bibitem{deneubourg1990}
J.-L.~Deneubourg, S.~Aron, S.~Goss, and J.-M.~Pasteels.
The self-organizing exploratory pattern of the Argentine ant.
\emph{Journal of Insect Behavior}, 3(2):159--168, 1990.

% --- Neural stigmergy ---

\bibitem{frey1997}
U.~Frey and R.~G.~M.~Morris.
Synaptic tagging and long-term potentiation.
\emph{Nature}, 385:533--536, 1997.

\bibitem{bhatt2000}
D.~K.~Bhatt, S.~Bhatt, and N.~J.~Abbott.
Four-dimensional neuronal signaling by nitric oxide: A computational
analysis.
\emph{Journal of Neuroscience}, 20(3):1199--1207, 2000.

\bibitem{philippides2005}
A.~Philippides, P.~Husbands, and M.~O'Shea.
From synaptically localized to volume transmission by nitric oxide.
\emph{Journal of Physiology}, 2005.

% --- Companion paper ---

\bibitem{stigmergic2024}
Stigmergic attention: Closing the feedback loop between pheromone
fields and neural attention.
\emph{Companion paper}, 2024.

% --- Transformers and attention variants ---

\bibitem{vaswani2017attention}
A.~Vaswani, N.~Shazeer, N.~Parmar, J.~Uszkoreit, L.~Jones,
A.~N.~Gomez, \L.~Kaiser, and I.~Polosukhin.
Attention is all you need.
In \emph{NeurIPS}, 2017.

\bibitem{press2022alibi}
O.~Press, N.~Smith, and M.~Lewis.
Train short, test long: Attention with linear biases enables input
length extrapolation.
In \emph{ICLR}, 2022.

\bibitem{su2024rope}
J.~Su, M.~Ahmed, Y.~Lu, S.~Pan, W.~Bo, and Y.~Liu.
RoFormer: Enhanced transformer with rotary position embedding.
\emph{Neurocomputing}, 568:127063, 2024.

\bibitem{sun2023retnet}
Y.~Sun, L.~Dong, S.~Huang, S.~Ma, Y.~Xia, J.~Xue, J.~Wang, and F.~Wei.
Retentive network: A successor to transformer for large language models.
\emph{arXiv:2307.08621}, 2023.

\bibitem{gu2023mamba}
A.~Gu and T.~Dao.
Mamba: Linear-time sequence modeling with selective state spaces.
\emph{arXiv:2312.00752}, 2023.

\bibitem{peng2023rwkv}
B.~Peng, E.~Alcaide, Q.~Anthony, et~al.
RWKV: Reinventing RNNs for the transformer era.
In \emph{EMNLP Findings}, 2023.

\bibitem{fan2021feedback}
A.~Fan, T.~Lavril, E.~Grave, A.~Joulin, and S.~Sukhbaatar.
Addressing some limitations of transformers with feedback memory.
\emph{arXiv:2002.09402}, 2021.

\bibitem{ramsauer2021hopfield}
H.~Ramsauer, B.~Sch\"{a}fl, J.~Lehner, et~al.
Hopfield networks is all you need.
In \emph{ICLR}, 2021.

\end{thebibliography}

\end{document}
